\begin{afterword}
\summaryhead

The author read about half of Adam Frank's \textsc{The Constant Fire} before a video dialogue with Frank. Frank holds that science and religion share the experience of the sacred, the feeling the author has watching a shuttle launch. The book quotes William James's definition of religion as the experiences of individuals ``in their solitude,'' and develops the theme that the sacred is private. The author asks what ``private'' could mean, and concludes that it is a trick of ``Dark Side Epistemology'': what is private cannot be criticized.

Religion distorts the sacred through five habits that grew up to shield its original lie: mysteriousness, faith, the retreat to pure experience, separation from the ordinary world, and privacy. Each has a cost. So religion should not be salvaged, not even as ``spirituality''; the wise start over. The capacity to admit being entirely wrong separates scientific from religious experience. The post ends with the sacred as it should be: real, knowable, shared, ``the sacred mundane.''

\responsehead

The post has two parts: a reply to Frank's book, and a case against keeping anything of religion, even as spirituality. The first has less support than it seems to. The second argues well against five habits and then draws a conclusion that this argument does not support.

The reply to Frank reports the book in the author's words, from about half of it. The only words quoted from Frank are ``unique experience of the sacred.'' The one full quotation is from James, and its context gives a different reason for stressing the individual. James offered his definition ``arbitrarily,'' as a working limit for lectures that would leave churches and theology aside. His ``solitude'' marks the individual as against the institution, and he treated personal religion as primary because the founders of churches drew their power from direct experience. The post reads the stress on privacy as a way to shield experience from criticism, and supports the reading with a line of defense it writes itself. Neither James nor Frank is shown using privacy that way.

The five habits are called ``relics'' of Dark Side Epistemology: each once protected religion's original lie. For each, the post names a cost, and those costs are arguments against the habits on their own terms. The claim about how the habits arose gets one historical sequence, under ``Faith,'' which repeats a passage of ``Religion's Claim to be Non-Disprovable'' nearly word for word.

The argument is weakest in its verdict on spirituality. The post names no spiritual practice and shows none of them keeping the five habits; the one book it discusses is faulted only for its stress on privacy. The step from the claim that the habits grew up to defend a lie, to the verdict that what is left is ``still not good for you,'' is made only through the image of a poisoned cup. The author's ``The Genetic Fallacy'' had endorsed doubting whatever grew from a discredited source, but added that one is ``not licensed to reject them outright as conclusions.'' The costs show that the habits do harm. They do not show that every form of spirituality has them.

The ending states the author's view clearly: an experience of the sacred that is believable, real, knowable, shared and continuous with ordinary things. Each clause answers one of the five habits.

In short: a clear statement of the author's view of the sacred, set against a book read in part and a quotation from James whose context gives another reason for stressing the individual, and closed with a verdict on all spirituality that examines no spiritual practice.
\end{afterword}
